\documentclass[10pt,a4paper]{amsart}
\usepackage[margin=19mm]{geometry}
\usepackage{amsmath,amssymb,mathtools}
\usepackage[scr=rsfs,cal=euler]{mathalfa}
\usepackage{xfrac}
\usepackage[unicode,hidelinks,bookmarks=false]{hyperref}
\newcommand{\ie}{i.e.\ }
\newcommand{\cf}{cf.\ }
\newcommand{\resp}{respectively\ }
\newcommand{\Subsection}{\subsection}
\providecommand{\nobreakdash}{\nobreak\hbox{-}}
\setlength{\emergencystretch}{2em}
\multlinegap0pt
\usepackage{amssymb}
\usepackage{mathtools}
\usepackage[shortlabels]{enumitem}
\usepackage{float}
\usepackage{xcolor}
\newtheorem{theorem}{Theorem}
\title{\textbf{Analysis of Chaos and Bifurcation in Nonlinear two-delay differential equation}}
\author{Pragati Dutta*, Sachin Bhalekar}
\date{Source: arXiv:2604.07924v2 (CC BY 4.0)}
\begin{document}
\maketitle

\noindent\emph{Source and licence.} \textbf{Analysis of Chaos and Bifurcation in Nonlinear two-delay differential equation}. Version arXiv:2604.07924v2 (CC BY 4.0), distributed by arXiv under the Creative Commons Attribution 4.0 International licence, http://creativecommons.org/licenses/by/4.0/, as recorded in the arXiv licence metadata for this exact version and shown on the abstract page https://arxiv.org/abs/2604.07924v2. Full licence notice: \texttt{licenses/arxiv-2604.07924-CC-BY-4.0.txt}.\par\smallskip

\noindent\textit{Editorial scope.} Counted unit: the article's theorem thm:fixed\_tau2\_sync (source0.tex lines 1568--1594). The article's complete original proof of it is delivered with it (source0.tex lines 1596--1826, one \texttt{proof} environment of 231 lines). No mathematics is written, repaired or re-derived: the statement and the proof are the article's own lines, in the article's own notation, and no retained line is substituted or re-flowed.  Retained uncounted as the source-local context the proof needs: lines 1409--1415; lines 1417--1424.  Nothing was omitted from any retained span, so the package carries no omission record.  No retained line contains a citation, so the wrapper carries no bibliography and no bibliography entry.  No retained line contains a figure environment, an image-inclusion command or drawing code, so no image is delivered and the package declares no asset dependency.  Editorial formatting only, in addition: an AMS \texttt{amsart} preamble that restates the article's own declarations and macros used by the retained lines, namely the environments \texttt{theorem} on separate counters, together with the standard packages those lines need.  The retained environments print in this document as Theorem 1 in order of appearance, whereas the article prints the same items with the numbering of its own full document.  The complete native source of the article as distributed in the arXiv e-print for this exact version is bundled unchanged as \texttt{sources/198124/source0.tex}.
\begin{equation}
	\dot{x}_1(t)
	= -\gamma x_1(t)
	+ k\sin\big(x_1(t - \tau_1)\big)
	- k e^{-\gamma \tau_2}\sin\big(x_1(t - \tau_1 - \tau_2)\big),
	\label{eq:ms}
\end{equation}

\begin{equation}
	\dot{x}_2(t)
	= -\gamma x_2(t)
	+ k\sin\big(x_2(t - \tau_1)\big)
	- k e^{-\gamma \tau_2}\sin\big(x_2(t - \tau_1 - \tau_2)\big)
	+ C(x_1,x_2),
	\label{eq:ss}
\end{equation}

\begin{theorem}[Less Conservative Fixed-$\tau_2$ Synchronization]
\label{thm:fixed_tau2_sync}
Assume that $\gamma\geq0$, fix $\tau_2\geq0$, and let $\tau_1\geq0$ be arbitrary. If the feedback gain $\delta$ satisfies
\begin{equation}
\gamma+\delta>|k|\bigl(1+e^{-\gamma\tau_2}\bigr),
\label{eq:less_conservative_sync}
\end{equation}
then the master and slave systems~\eqref{eq:ms}--\eqref{eq:ss} are globally exponentially synchronized. More precisely, let $\eta>0$ be the unique positive solution of
\begin{equation}
\eta
=
(\gamma+\delta)
-|k|e^{-\eta\tau_1}
-|k|e^{-\eta(\tau_1+\tau_2)}.
\label{eq:fixed_tau2_rate}
\end{equation}
Then, for every $\sigma\in(0,\eta)$, there exists a constant $M_\sigma>0$, determined by the initial error history, such that
\begin{equation}
|e(t)|\leq M_\sigma e^{-\sigma t},
\qquad t\geq0.
\label{eq:fixed_tau2_decay}
\end{equation}
Consequently,
\[
\lim_{t\to\infty}e(t)=0.
\]
\end{theorem}

\begin{proof}
Let
\[
v(t)=|e(t)|.
\]
Since $v(t)$ may fail to be differentiable at points where $e(t)=0$, we work with its upper right Dini derivative
\[
D^+v(t)
=
\limsup_{h\to0^+}
\frac{v(t+h)-v(t)}{h}.
\]
Using the error equation and the inequality
\[
|\sin u-\sin v|\leq |u-v|,
\]
we obtain
\begin{equation}
D^+v(t)
\leq
-(\gamma+\delta)v(t)
+|k|v(t-\tau_1)
+|k|e^{-\gamma\tau_2}
v(t-\tau_1-\tau_2),
\qquad t\geq0.
\label{eq:fixed_tau2_dini}
\end{equation}
This is a two-delay differential inequality of Halanay type.

Define
\[
F(r)
=
\gamma+\delta-r
-|k|e^{r\tau_1}
-|k|e^{-\gamma\tau_2+r(\tau_1+\tau_2)},
\qquad r\geq0.
\]
Then
\[
F(0)
=
\gamma+\delta-|k|-|k|e^{-\gamma\tau_2}>0
\]
by~\eqref{eq:less_conservative_sync}, whereas
\[
F(r)\to-\infty
\qquad\text{as }r\to\infty.
\]
Furthermore,
\[
F'(r)
=
-1
-|k|\tau_1e^{r\tau_1}
-|k|e^{-\gamma\tau_2}
(\tau_1+\tau_2)e^{r(\tau_1+\tau_2)}
<0.
\]
Thus, $F$ is strictly decreasing and has a unique positive zero $\eta$. The equation $F(\eta)=0$ is precisely~\eqref{eq:fixed_tau2_rate}.

Fix any $\sigma\in(0,\eta)$. Since $F$ is strictly decreasing,
\[
F(\sigma)
=
\gamma+\delta-\sigma
-|k|e^{\sigma\tau_1}
-|k|e^{-\gamma\tau_2}
e^{\sigma(\tau_1+\tau_2)}
>0.
\label{eq:positive_margin_sigma}
\]

Let
\[
e_0(s)=e(s),
\qquad
s\in[-(\tau_1+\tau_2),0],
\]
denote the initial error history and set
\[
M_\sigma
=
\max_{-(\tau_1+\tau_2)\leq s\leq0}
e^{\sigma s}|e_0(s)|.
\]
If $M_\sigma=0$, then the two initial histories coincide, and uniqueness of solutions gives
\[
e(t)\equiv0.
\]
Suppose therefore that $M_\sigma>0$ and define
\[
w(t)=M_\sigma e^{-\sigma t},
\qquad
t\geq-(\tau_1+\tau_2).
\]
By the definition of $M_\sigma$, for every
\[
s\in[-(\tau_1+\tau_2),0],
\]
we have
\[
e^{\sigma s}|e_0(s)|
\leq M_\sigma.
\]
Hence
\[
v(s)=|e_0(s)|
\leq
M_\sigma e^{-\sigma s}
=
w(s),
\]
and therefore
\[
v(s)\leq w(s),
\qquad
-(\tau_1+\tau_2)\leq s\leq0.
\]
Thus, the comparison function $w$ dominates the complete initial error history.

We now prove that
\[
v(t)\leq w(t),
\qquad t\geq0.
\]
Suppose, to the contrary, that there exists $t>0$ such that
\[
v(t)>w(t).
\]
By continuity of $v$ and $w$ and the fact that $v(s)\leq w(s)$ on the initial interval, there exists a first contact time $t_*>0$ such that
\[
v(t_*)=w(t_*),
\qquad
v(t)\leq w(t)
\quad
\text{for }
-(\tau_1+\tau_2)\leq t\leq t_*,
\]
and $v(t)>w(t)$ for times immediately to the right of $t_*$.

Consequently,
\[
D^+(v-w)(t_*)\geq0.
\]
On the other hand, using~\eqref{eq:fixed_tau2_dini} together with
\[
v(t_*-\tau_1)\leq w(t_*-\tau_1)
\]
and
\[
v(t_*-\tau_1-\tau_2)
\leq
w(t_*-\tau_1-\tau_2),
\]
we obtain
\begin{align*}
D^+v(t_*)
&\leq
-(\gamma+\delta)w(t_*)
+|k|w(t_*-\tau_1)\\
&\quad
+|k|e^{-\gamma\tau_2}
w(t_*-\tau_1-\tau_2).
\end{align*}
Since
\[
w(t)=M_\sigma e^{-\sigma t},
\]
we have
\[
w(t_*-\tau_1)
=
e^{\sigma\tau_1}w(t_*),
\]
and
\[
w(t_*-\tau_1-\tau_2)
=
e^{\sigma(\tau_1+\tau_2)}w(t_*).
\]
Therefore,
\begin{align*}
D^+v(t_*)
&\leq
\left[
-(\gamma+\delta)
+|k|e^{\sigma\tau_1}
+|k|e^{-\gamma\tau_2}
e^{\sigma(\tau_1+\tau_2)}
\right]w(t_*)\\
&=
-\left[\sigma+F(\sigma)\right]w(t_*)\\
&<
-\sigma w(t_*).
\end{align*}
Since
\[
w'(t_*)=-\sigma w(t_*),
\]
and $w$ is continuously differentiable, we obtain
\[
D^+(v-w)(t_*)
=
D^+v(t_*)-w'(t_*)
<0,
\]
which contradicts
\[
D^+(v-w)(t_*)\geq0.
\]
Hence
\[
v(t)\leq w(t),
\qquad t\geq0.
\]
Consequently,
\[
|e(t)|
=
v(t)
\leq
M_\sigma e^{-\sigma t},
\qquad t\geq0,
\]
which proves~\eqref{eq:fixed_tau2_decay}. Since $\sigma>0$,
\[
\lim_{t\to\infty}|e(t)|=0.
\]
Thus, complete synchronization is achieved.
\end{proof}

\end{document}
